\addcontentsline{toc}{subsubsection}{Solution 2. IMO 2013 -- Problem 5}

\begin{dlsummary}[title={\bfseries Solution 2. IMO 2013 -- Problem 5},breakable]

\begin{proof}
From $f(a)f(1)\geq f(a)$, we have $f(1)\geq1$, hence
$f(n)\geq n$ for positive integers $n$. Also,
$f(q)f(p/q)\geq f(p)>0$ shows $f(r)>0$ for every positive rational
$r$. Thus superadditivity makes $f$ strictly increasing.

Submultiplicativity gives $f(a^n)\leq a^n$, while
$f(a^n)f(a^{-n})\geq f(1)\geq1$ yields $f(a^{-n})\geq a^{-n}$.
For rational $x>0$, set $m_n=\lfloor a^n x\rfloor$. For large $n$,
monotonicity and superadditivity give
\[
f(x)\geq f(m_na^{-n})\geq m_nf(a^{-n})
\geq\frac{\lfloor a^n x\rfloor}{a^n}\longrightarrow x.
\]
Hence $f(x)\geq x$.

Applying this to $a^n$ shows $f(a^n)=a^n$. For any positive
rational $x$, choose $n$ with $a^n>x$. Then
\[
a^n=f(a^n)\geq f(x)+f(a^n-x)\geq f(x)+a^n-x,
\]
so $f(x)\leq x$. Therefore $f(x)=x$ for all positive rationals;
the identity satisfies both inequalities and the fixed-point condition.
\end{proof}

\end{dlsummary}